
\lussec 10. Concluding remarks

Through its extension to the quark fields,
the gradient flow becomes a powerful tool for 
studies of the spontaneous breaking of chiral symmetry in QCD
and the physics of the light pseudo-scalar mesons.
The focus in this paper has been on
the theoretical framework and its implementation in 
one of the widely used formulations of lattice QCD. 
In view of the renormalizability properties of the flow,
the choice of the lattice regularization is however expected to be
irrelevant in the continuum limit
as long as the locality and gauge invariance 
of the lattice theory are guaranteed.
The theoretical analysis obviously applies to 
field theories
with other gauge groups and fermion multiplets as well.

In the cases worked out so far,
the application of the gradient flow in numerical lattice QCD
offers important technical advantages 
with respect to the established computational strategies. 
Apart from the ones already mentioned in sect.~1, 
there are potentially many further uses of the flow.
The flavour singlet channel, for example,
has not been considered yet and the time-dependent condensate
is likely to be an easily accessible chiral
order parameter in QCD at non-zero temperatures.
Moreover, the application of the 
flow in QCD-like theories near the conformal window may 
conceivably lead to interesting qualitative insights.

